\section*{الفصل \ref{ch:b2:neurons-synapses} --- العصبونات وكمونات الفعل والمشابك}

\begin{solution}{exo:b2:neurons-synapses:1}
التغصنات تتلقى المشابك؛ وجسم الخلية يحمل النواة
ويكامل؛ \omterm{prop:b2:neurons-synapses:integration}{وربوة المحوار} تقرر؛ \omterm{def:b2:neurons-synapses:neuron}{والمحوار} يوصّل؛
والنهايات تطلق الناقل. \omterm{def:b2:neurons-synapses:neuron}{والميالين} غشاء ملفوف من
\omterm{def:b2:neurons-synapses:neuron}{خلية دبقية}، يكاد يخلو من الهيولى؛ وهو يعزل \omterm{def:b2:neurons-synapses:neuron}{المحوار} فتقفز
الدفعة بين العقد وتسير أسرع مئة مرة
بكلفة أقل.
\end{solution}

\begin{solution}{exo:b2:neurons-synapses:2}
زوال استقطاب إلى العتبة؛ ثم الطور الصاعد --- تنفتح قنوات
الصوديوم فولطية البوابة، وتغذية راجعة موجبة؛ ثم الذروة قرب $E_{\mathrm{Na}}$ ---
تُعطَّل قنوات الصوديوم؛ ثم الطور الهابط --- تنفتح قنوات البوتاسيوم
فولطية البوابة؛ ثم الهبوط دون الراحة --- قنوات البوتاسيوم ما تزال مفتوحة، قرب
$E_{\mathrm K}$؛ ثم التعافي إذ تنغلق. وهو لا يقبل التدرج لأن
التغذية الراجعة الموجبة إما أن تمضي إلى النهاية وإما ألا تكون؛ وهو أحادي الاتجاه
لأن الغشاء الذي قُطع لتوّه حَرِن (فقنوات الصوديوم
معطَّلة).
\end{solution}

\begin{solution}{exo:b2:neurons-synapses:3}
تبلغ الدفعة النهاية؛ فتنفتح قنوات الكالسيوم فولطية البوابة؛
ويطلق الكالسيوم اندماج الحويصلات؛ وينتشر الناقل عبر
الشق؛ ويرتبط بالمستقبلات التالية \omterm{def:b2:neurons-synapses:synapse}{للمشبك}؛ فتنفتح القنوات (أو يقلب \omterm{prop:b2:cell-signalling:gpcr}{بروتين
G})؛ \omterm{def:b2:neurons-synapses:synapse}{فكمون تالٍ للمشبك}؛ ثم يُزال الناقل بإنزيم أو
بناقل؛ ويُستعاد غشاء الحويصل.
\end{solution}

\begin{solution}{exo:b2:neurons-synapses:4}
الكهربائي: بلا تأخر، وفي الاتجاهين عادةً، وبلا تضخيم، وبلا
تثبيط (إذ لا يمر غير التيار). والكيميائي: نصف مليثانية،
وأحادي الاتجاه، وبتضخيم (فالدفعة الواحدة تطلق مئات الكمومات)،
وبتثبيط ممكن (قنوات للكلور أو البوتاسيوم)، وقابل
للتعديل.
\end{solution}

\begin{solution}{exo:b2:neurons-synapses:5}
% Long unbreakable runs (a chain of numerals, a Latin binomial) are one
% directional box under bidi=basic and cannot be broken; \sloppy must be in
% force when the paragraph is BROKEN, hence the \par before \endgroup.
\begingroup\sloppy
من $E = (61.5/z)\log_{10}(c_{\text{خارج}}/c_{\text{داخل}})$:
البوتاسيوم \qty{-89}{mV}؛ والصوديوم \qty{+61}{mV}؛ والكلور، عند $z =
-1$، $-61.5\log_{10} 11 = \qty{-64}{mV}$؛ والكالسيوم، عند $z = 2$،
$30.75\times 4.3 = \qty{+132}{mV}$. وعند \qty{-70}{mV}:
يخرج البوتاسيوم، ويدخل الصوديوم، ويخرج الكلور قليلًا (فالغشاء
دون توازنه)، ويدخل الكالسيوم بقوة.
\par\endgroup
\end{solution}

\begin{solution}{exo:b2:neurons-synapses:6}
عند $20 : 1$: $(20\times(-89) + 61)/21 = \qty{-82}{mV}$، وهي حالة
الراحة قرب $E_{\mathrm K}$. وعند $1 : 20$: $(-89 + 20\times 61)/21 =
\qty{+54}{mV}$، وهي الذروة، قرب $E_{\mathrm{Na}}$. وعند $50 : 1$: $(50\times
(-89) + 61)/51 = \qty{-86}{mV}$، وهو الهبوط دون الراحة مع قنوات بوتاسيوم
إضافية مفتوحة.
\end{solution}

\begin{solution}{exo:b2:neurons-synapses:7}
$\sqrt{500} = 22$: أي \qty{22}{m/s}. ولبلوغ \qty{100}{m/s} كان على
\omterm{def:b2:neurons-synapses:neuron}{محوار} بلا \omterm{def:b2:neurons-synapses:neuron}{ميالين} أن يكون $d = 10^{4}$ \unit{\micro m}، أي
سنتيمترًا. والمتر الواحد يستغرق \qty{10}{ms} في الليف ذي \omterm{def:b2:neurons-synapses:neuron}{الميالين}،
و\qty{45}{ms} في \omterm{def:b2:neurons-synapses:neuron}{محوار} الحبار، وثانية كاملة في ليف قطره \qty{1}{\micro m}.
\end{solution}

\begin{solution}{exo:b2:neurons-synapses:8}
$m = \ln(200/74) = 1.0$؛ و$P(1) = e^{-1} = 0.37$، و$P(2) = 0.18$، و$P(3)
= 0.06$. والاستجابة المتوسطة $1.0\times 0.5 = \qty{0.5}{mV}$.
\end{solution}

\begin{solution}{exo:b2:neurons-synapses:9}
$1/\qty{2}{ms} = 500$ دفعة في الثانية. وشدة اللمس
يرمّزها تواتر الدفعات (وعدد العصبونات
المجنَّدة)؛ أما الدفعات نفسها فمتماثلة.
\end{solution}

\begin{solution}{exo:b2:neurons-synapses:10}
$\lambda = \sqrt{R_m d/4R_i}$: أي \qty{0.5}{mm} و\qty{1.6}{mm}
و\qty{11}{mm}. ومع \omterm{def:b2:neurons-synapses:neuron}{الميالين}، عند \qty{10}{\micro m}: $1.6\times\sqrt{200} =
\qty{22}{mm}$. والقطعة بين عقدتين طولها \qty{1}{mm} لا تفقد غير $1 - e^{-1/22}
= \qty{4}{\%}$ من الإشارة، فتُبلغ العقدة التالية العتبةَ
بيسر. وعبر \qty{3}{mm} منزوعة \omterm{def:b2:neurons-synapses:neuron}{الميالين} تهبط الإشارة إلى
$e^{-3/1.6} = \qty{15}{\%}$ --- أي من دفعة قدرها \qty{100}{mV} نحو
عتبة \qty{15}{mV}: فيخفق التوصيل أو يصير غير موثوق.
\end{solution}

\begin{solution}{exo:b2:neurons-synapses:11}
التترودوتوكسين: لا تيار صوديوم، ولا دفعة، ولا نقل.
ورباعي إيثيل الأمونيوم: لا تيار بوتاسيوم، ودفعة مطوَّلة، وناقل
أكثر يُطلق في الدفعة. وبلا كالسيوم: دفعات سوية،
وبلا إطلاق، وبلا نقل. والكورار: إطلاق سوي، ومستقبلات
محصورة، وبلا كمون صفيحة انتهائية --- أي شلل. ومثبط
الكولين إستيراز: يدوم الأستيل كولين، وتبقى المستقبلات مفتوحة، فيزول استقطاب
العضلة ثم لا تستطيع إعادة استقطابها --- أي شلل بالفائض.
وذيفان الوشيقية: لا تستطيع الحويصلات الاندماج، فلا إطلاق --- أي شلل.
\end{solution}

\begin{solution}{exo:b2:neurons-synapses:12}
يشتري التأخر: نقلًا في اتجاه واحد فقط؛ وتضخيمًا
(فالدفعة الواحدة تطلق مئات الكمومات وتستطيع قيادة
خلية أكبر)؛ وتثبيطًا، وهو ما لا يقدر عليه وصل كهربائي؛
ولدونةً، إذ يستطيع عدد الكمومات والمستقبلات أن يتغير
بالاستعمال، وهو كيف تتعلم الدارات. أما المشابك الكهربائية فتُستعمل
حيث يهم التزامن والسرعة أكثر من الحساب:
دارات الهرب، وعضلة \omterm{def:b2:heart:anatomy}{القلب}، وبعض الشبكات المثبطة.
\end{solution}

\begin{solution}{pb:b2:neurons-synapses:1}
\textbf{1.} $E_{\mathrm K} = 61.5\log_{10}(4/140) = \qty{-95}{mV}$؛
و$E_{\mathrm{Na}} = 61.5\log_{10}(145/12) = \qty{+67}{mV}$.
\textbf{2.} $(25\times(-95) + 67)/26 = \qty{-89}{mV}$.
\textbf{3.} $E_{\mathrm K} = 61.5\log_{10}(8/140) = \qty{-76}{mV}$؛
و$V = (25\times(-76) + 67)/26 = \qty{-71}{mV}$: أي أقرب إلى العتبة،
فأشد استثارية في البداية --- وأقل استثارية إن دام، إذ تُعطَّل قنوات
الصوديوم عند المستوى الزائل الاستقطاب.
\textbf{4.} تنفد التدرجات على مدى ساعات إذ يرشح الصوديوم داخلًا
والبوتاسيوم خارجًا؛ وينجرف الكمون نحو الصفر؛ ومع توقف
المضخة تجذب شوارد الخلية السالبة غير النفوذة الكلورَ والماء،
فتنتفخ.
\textbf{5.} السطح في المليمتر $\pi\times 15\times 10^{-6}\times
10^{-3} = \qty{4.7e-8}{m^2}$؛ و$C = \qty{4.7e-10}{F}$؛ و$Q = CV =
\qty{4.7e-11}{C}$: أي $2.9\times 10^{8}$ شاردة.
\textbf{6.} الحجم في المليمتر $\pi(7.5\times 10^{-6})^{2}\times
10^{-3} = \qty{1.8e-13}{m^3} = \qty{1.8e-10}{L}$، ويحمل $12\times
10^{-3}\times 1.8\times 10^{-10} = 2.1\times 10^{-12}$ مول، أي $1.3\times
10^{12}$ شاردة: فالتغير $2\times 10^{-4}$.
\textbf{7.} $2.9\times 10^{8}/3 \approx 10^{8}$ من ATP في المليمتر لكل
دفعة.
\textbf{8.} $0.9/90 = \qty{10}{ms}$ في كل اتجاه.
\textbf{9.} عاريًا: $\sqrt{1\times 15\times 10^{-6}/4} = \qty{1.9}{mm}$؛
وبميالين: $\times\sqrt{250} = \qty{31}{mm}$.
\textbf{10.} عاريًا: $e^{-1.5/1.9} = 0.45$، أي \qty{45}{mV}؛ وبميالين:
$e^{-1.5/31} = 0.95$، أي \qty{95}{mV}. وكلاهما يتجاوز عتبة \qty{15}{mV}؛
لكن على \omterm{def:b2:neurons-synapses:neuron}{المحوار} العاري أن يشحن غشاءه كله على امتداد
الطريق، وهو ما يجعله بطيئًا، في حين لا يكاد يُفقد أو يُشحن شيء
تحت \omterm{def:b2:neurons-synapses:neuron}{الميالين}.
\textbf{11.} $\sqrt{15} = \qty{3.9}{m/s}$: أي \qty{0.23}{s} في كل اتجاه،
ونحو نصف ثانية للرحلة كلها --- وهو أبطأ من أن يكفي منعكسًا
عليه أن يتدارك عثرة.
\textbf{12.} $e^{-4/1.9} = 0.12$: أي \qty{12}{mV} عند العقدة البعيدة، دون
العتبة --- فتُحصر الدفعة ويضيع المنعكس.
\textbf{13.} خلف الدفعة تكون قنوات الصوديوم معطَّلة
مليثانية أو اثنتين، فلا يمكن استثارة الغشاء الذي قُطع لتوّه
ولا تستطيع الموجة إلا المضي قُدُمًا.
\textbf{14.} $m = \ln(300/111) = 1.0$.
\textbf{15.} $P(1) = 0.37$، و$P(2) = 0.18$، و$P(3) = 0.06$: أي نحو 110 و55
و18 من المحاولات 300.
\textbf{16.} $1.0\times 0.4 = \qty{0.4}{mV}$.
\textbf{17.} $P(0) = e^{-60} \approx 10^{-26}$: أي أبدًا؛ والاستجابة المتوسطة
\qty{24}{mV}، فوق عتبة \qty{15}{mV}: \omterm{def:b2:neurons-synapses:synapse}{فمشبك} واحد كهذا
يطلق \omterm{def:b2:neurons-synapses:neuron}{العصبون} الحركي.
\textbf{18.} عشرون مشبكًا فعالة معًا تجمع تياراتها
(تجميعًا مكانيًا)، جمعًا دون الخطي إذ يقترب الكمون من
كمون انعكاس القنوات؛ وحتى عند $m = 1$ لكل منها كانت
تعطي \qty{8}{mV}، وعند $m$ السوي تعطي أكثر من العتبة بكثير.
\textbf{19.} فتح قنوات الكلور عند $E_{\mathrm{Cl}} = V_{\text{الراحة}}$
يضيف ناقلية دون تحريك الكمون؛ فيجري التيار
المنبّه الآن عبر غشاء أشد رشحًا فينتج زوال استقطاب
أصغر --- وهو التثبيط بالتحويل.
\textbf{20.} $10 + 0.6 + 10 + 0.6 + 5 = \qty{26}{ms}$، في مقابل
\qty{30}{ms} ملحوظة.
\textbf{21.} تفعيل المستقبل نفسه في المغزل العضلي،
وانتشار كمون فعل الليف العضلي على امتداد الليف
(أمتار في الثانية على مدى سنتيمترات)، والكمون الزمني الميكانيكي
قبل ظهور القوة.
\textbf{22.} كل \omterm{def:b2:neurons-synapses:synapse}{مشبك} يضيف نصف مليثانية وكل \omterm{def:b2:neurons-synapses:neuron}{عصبون} بيني
يضيف خلية؛ لكن العصبونات البينية تتيح \omterm{def:b2:heart:anatomy}{قلب} الإشارة
(بتثبيط المضاد)، ودمجها بغيرها، وتعديلها من
الدماغ --- أي منعكسًا يمكن تشكيله.
\textbf{23.} $90 - 6.2 = \qty{84}{ms}$ من التوصيل لمسافة \qty{1.8}{m}:
أي نحو \qty{21}{m/s}.
\textbf{24.} نصف التيار الداخل عند كل فولطية: فترتفع العتبة،
وتصير الدفعة أصغر وأبطأ، ويهبط عامل الأمان عند كل
عقدة، وقد تخفق الدفعة؛ فيضعف المنعكس أو
يزول.
\textbf{25.} \omterm{thm:b2:neurons-synapses:chord}{كمون الراحة} \qty{-89}{mV}؛ و\qty{10}{ms} في كل اتجاه؛
و$m = 1.0$؛ والكمون الزمني المحسوب \qty{26}{ms}.
\end{solution}
